\documentclass[../main.tex]{subfiles}
\begin{document}

\chapter{Branches}
\lessoninfo{
  video={Lesson 4, videos 1--48},
  textbook={H\&P \cite{hennessy2017} §3.3, §3.9, §C.2},
  keywords={branch prediction, misprediction penalty, predict not taken, branch target buffer, branch history table, two-bit predictor, history-based predictor, PShare, GShare, tournament predictor, hierarchical predictor, return address stack, predecoding},
}

Lesson 3 showed that control dependences are hazards in every pipeline: the
processor has to fetch the instructions that follow a branch long before it
knows where the branch goes, and every wrong choice flushes work.  This lesson
studies how the fetch stage makes that choice.  Starting from the simplest
policy---always continue with the next instruction---it develops the
structures used in real processors: branch target buffers, saturating
counters, predictors that learn patterns from branch history, combinations of
several predictors, and a dedicated predictor for function returns.  Along the
way it shows why prediction accuracy matters more the deeper the pipeline.

\section{Branches in a pipeline}
\label{sec:l04-branch-pipeline}

Consider a conditional branch \texttt{BEQ R1,R2,L} in the five-stage pipeline
of Lesson 3.\source{Lesson 4, videos 1--3; H\&P §C.2.}  If R1 equals R2,
execution continues at label~L; the branch is then a \term{taken
branch}[成立分岐], and the address of~L is its \term{branch target
address}[分岐先アドレス].

Otherwise execution continues with the next instruction in memory, at PC+4,
and the branch is a \term{not-taken branch}[不成立分岐].  The pipeline learns
which case applies only when the branch has passed through three stages: IF
fetches it, ID reads R1 and R2, and EX compares them and computes the target.
The branch is \emph{resolved} at the end of EX (\cref{fig:l04-branch-pipeline}).

The fetch stage cannot wait for that.  In cycle~2, while the branch is in ID,
and again in cycle~3, it must load \emph{some} instruction.  When the branch
resolves, two instructions are already in the pipeline behind it.  If they
are the instructions that really follow the branch, nothing is lost; if not,
they are flushed and two cycles are wasted.  Fetching nothing during these
cycles would waste the same two cycles unconditionally, so it always pays to
fetch the instruction that is most likely to be the right one.

\begin{figure}[htbp]
  \centering
  % A branch in the five-stage pipeline: the two instructions fetched while the
% branch is in ID and EX are predictions, checked when the branch resolves.
\begin{tikzpicture}[x=0.95cm,y=0.55cm, font=\footnotesize\sffamily,
  st/.style={draw, minimum width=0.9cm, minimum height=0.5cm, inner sep=0pt, fill=ln-box},
  pr/.style={st, fill=ln-accent!15}]
  \foreach \c in {1,...,6} { \node at (\c,0.8) {\c}; }
  \node[anchor=east] at (0.5,0.8) {\iflangja{サイクル}{cycle}};
  \node[anchor=east] at (0.5,0) {\texttt{BEQ R1,R2,L}};
  \foreach \s/\lab in {0/IF,1/ID,2/EX,3/MEM,4/WB} { \node[st] at (\s+1,0) {\lab}; }
  \node[anchor=east] at (0.5,-1) {\iflangja{予測した命令 1}{predicted instr.\ 1}};
  \node[pr] at (2,-1) {IF}; \node[pr] at (3,-1) {ID};
  \node[anchor=east] at (0.5,-2) {\iflangja{予測した命令 2}{predicted instr.\ 2}};
  \node[pr] at (3,-2) {IF};
  \draw[ln-caution, thick, dashed] (3.5,1.15) -- (3.5,-2.6)
    node[below, font=\scriptsize] {\iflangja{分岐解決（EX の終わり）}{branch resolved (end of EX)}};
  \node[anchor=west, align=left, font=\scriptsize, text=ln-muted] at (3.75,-1.5)
    {\iflangja{予測が正しい: 両命令とも続行\\予測ミス: 両命令をフラッシュ}{correct prediction: both continue\\misprediction: both are flushed}};
\end{tikzpicture}

  \caption{A branch in the five-stage pipeline.  The two instructions fetched
    before the branch resolves are predictions; they are flushed if the
    prediction was wrong.}
  \label{fig:l04-branch-pipeline}
\end{figure}

\begin{definition}[Branch prediction]
  \label{def:l04-branch-prediction}
  The choice, made in the fetch stage, of the address from which the next
  instruction is fetched, before the current instruction has been decoded
  or executed, is called \term{branch prediction}[分岐予測].
\end{definition}

\begin{definition}[Misprediction]
  A \term{branch misprediction}[分岐予測ミス] occurs when the branch, once
  resolved, turns out to go elsewhere than predicted.  The instructions
  fetched from the wrong address are flushed and fetching restarts at the
  correct address.
\end{definition}

\subsection{What the predictor knows}

At the time of the fetch the only available information is the PC of the
instruction being fetched; the instruction itself has not even been decoded
yet.  From this address alone the predictor has to answer three questions:
\begin{enumerate}
  \item Is the instruction a branch at all?
  \item If it is, is it taken?
  \item If it is taken, what is its target address?
\end{enumerate}
The answer is a single number, the next PC: PC+4 for a non-branch or a
not-taken branch, the target address for a taken
branch.\margin{Unconditional jumps, calls and returns are branches that are
always taken; for them only questions 1 and~3 matter.}

\section{Prediction accuracy and performance}
\label{sec:l04-accuracy}

With a predictor in place, the only cycles lost to control hazards are those
spent on mispredicted paths.\source{Lesson 4, videos 4--5; H\&P §C.2.}
Specializing \cref{eq:pipeline-cpi} to branches gives
\begin{equation}
  \text{CPI} \;=\; 1 \;+\;
  \frac{\text{mispredictions}}{\text{instruction}} \times
  \frac{\text{penalty}}{\text{misprediction}} .
  \label{eq:l04-cpi}
\end{equation}

\begin{definition}[Accuracy and penalty]
  The \term{branch prediction accuracy}[分岐予測精度] is the fraction of
  branches whose outcome is predicted correctly, so that
  \begin{equation*}
    \frac{\text{mispredictions}}{\text{instruction}} =
    \frac{\text{branches}}{\text{instruction}} \times
    (1 - \text{accuracy}) ,
  \end{equation*}
  assuming that non-branch instructions are always predicted correctly, with
  PC+4 as the next PC.  The \term{misprediction penalty}[予測ミスペナルティ] is
  the number of cycles lost per misprediction; in a pipeline that holds one
  instruction per stage it equals the number of instructions flushed.  A
  branch resolved in stage $r$ has a penalty of $r-1$ cycles.
\end{definition}

The two factors of \cref{eq:l04-cpi} have different origins.  The first is
set by the predictor (and by how many branches the program executes); the
second is set by the pipeline, namely by how late branches are resolved.  A
better predictor shrinks the first factor, while a deeper pipeline enlarges
the second, so the same predictor costs more in a deeper pipeline.  A perfect
predictor would reach the ideal CPI of~1 regardless of the depth.

\begin{example}
  \label{ex:l04-cpi}
  In a program \SI{15}{\percent} of the instructions are branches and the
  predictor is \SI{92}{\percent} accurate.  Mispredictions per instruction
  are $0.15 \times 0.08 = 0.012$.  If branches are resolved in EX (stage~3,
  penalty~2), $\text{CPI} = 1 + 0.012 \times 2 = 1.024$.  In a pipeline that
  resolves branches in stage~11 (penalty~10), the same predictor gives
  $\text{CPI} = 1 + 0.012 \times 10 = 1.12$.
\end{example}

\section{Predict not taken}
\label{sec:l04-not-taken}

The alternative to predicting is to refuse to predict: the fetch stage simply
waits until it knows the next PC.\source{Lesson 4, videos 6--8; H\&P §C.2.}
Since even whether an instruction is a branch is known only after ID, the
wait affects every instruction.  A non-branch is followed by a fetch two
cycles after its own fetch, and a branch, taken or not, by a fetch three
cycles later.

\begin{table}[b]
  \centering\small
  \caption{Cycles from the fetch of an instruction until the correct next
    instruction is fetched, in the five-stage pipeline.}
  \label{tab:l04-not-taken}
  \begin{tabular}{@{}lcc@{}}
    \toprule
    Instruction        & Refuse to predict & Predict not taken \\
    \midrule
    Taken branch       & 3 & 3 \\
    Not-taken branch   & 3 & 1 \\
    Not a branch       & 2 & 1 \\
    \bottomrule
  \end{tabular}
\end{table}

The simplest real predictor is to \term{predict not taken}[不成立予測]:
always fetch from PC+4.  This needs no memory and no hardware beyond what the
fetch stage already has, since it increments the PC anyway.  It is right for
every non-branch and every not-taken branch, and wrong only for taken
branches, which cost three cycles as before.  \Cref{tab:l04-not-taken}
compares the two policies; predict not taken is never worse.

Its accuracy, however, is poor.  Most branches are taken---loop branches
jump back most of the time, and unconditional jumps, calls and returns are
always taken---so predict not taken mispredicts the majority of branches.
With the typical statistics used in Lesson 3 it mispredicts on
\SI{12}{\percent} of all instructions, giving the CPI of \cref{eq:branch-cpi}.

\begin{note}[Several predictions in flight]
  In a deep pipeline several branches can be between IF and the resolution
  stage at the same time, each carrying the prediction made when it was
  fetched.  When a branch resolves as mispredicted, everything fetched after
  it is flushed, including later branches together with their predictions.
  A branch fetched on a wrong path therefore never causes a flush of its own:
  only mispredictions of branches on the correct path cost cycles.
\end{note}

\section{The need for better prediction}
\label{sec:l04-better}

How much a better predictor is worth depends on the
pipeline.\source{Lesson 4, videos 9--11; H\&P §3.3.}  A processor that
handles $w$ instructions per stage flushes $w$ times as many instructions per
misprediction as a one-instruction-wide pipeline of the same depth, so the
execution time relative to perfect prediction is
\begin{equation}
  \frac{T}{T_{\text{perfect}}} \;=\; 1 +
  \frac{\text{mispredictions}}{\text{instruction}} \times
  (\text{instructions flushed per misprediction}) ,
  \label{eq:l04-relative}
\end{equation}
which for $w=1$ is exactly the CPI of \cref{eq:l04-cpi}.\margin{Real
numbers: a Skylake-class core loses about 17 cycles per misprediction, with a
224-entry reorder buffer to refill (Intel optimization manual).}

\begin{example}
  \label{ex:l04-better}
  Let \SI{15}{\percent} of the instructions be branches, \SI{70}{\percent} of
  them taken, so that predict not taken is \SI{30}{\percent} accurate.
  \Cref{tab:l04-better} evaluates \cref{eq:l04-relative} for three
  predictors and three pipelines: a five-stage pipeline (2 instructions
  flushed), a pipeline resolving branches in stage~13 (12 flushed), and the
  same depth with 3 instructions per stage (36
  flushed).\caution{Mispredictions are counted per \emph{instruction}, not per
  branch: multiply $1-\text{accuracy}$ by the fraction of instructions that
  are branches.}
\end{example}

\begin{table}[htbp]
  \centering\small
  \caption{Execution time relative to perfect prediction for the branch
    statistics of \cref{ex:l04-better}.}
  \label{tab:l04-better}
  \begin{tabular}{@{}lcccc@{}}
    \toprule
    Predictor & Mispred./instr. & Stage 3 & Stage 13 & Stage 13, 3 wide \\
    \midrule
    Not taken (\SI{30}{\percent}) & 0.105  & 1.21  & 2.26  & 4.78 \\
    \SI{90}{\percent} accurate    & 0.015  & 1.03  & 1.18  & 1.54 \\
    \SI{99}{\percent} accurate    & 0.0015 & 1.003 & 1.018 & 1.054 \\
    \bottomrule
  \end{tabular}
\end{table}

Three observations follow.\tip{(mispredictions per instruction) $\times$
(instructions flushed) is the slowdown itself: $0.015 \times 12 = 0.18$, the
\SI{18}{\percent} of the middle row.}  In the shallow pipeline even predict not taken
costs only \SI{21}{\percent}, and improving from \SI{90}{\percent} to
\SI{99}{\percent} accuracy gains less than \SI{3}{\percent}.  In the deep
pipeline predict not taken more than doubles the execution time, and the
same improvement is worth a speedup of $1.18/1.018 \approx 1.16$.  In the deep
and wide pipeline it is worth $1.54/1.054 \approx 1.46$.  Modern processors
are both deep and wide, which is why they invest heavily in branch prediction.

\begin{keypoint}
  What matters is the \emph{misprediction} rate, not the accuracy.  Going
  from \SI{90}{\percent} to \SI{99}{\percent} accuracy looks like a small
  step but removes \SI{90}{\percent} of the mispredictions, and the deeper and
  wider the pipeline, the more each remaining misprediction costs.
\end{keypoint}

\section{The branch target buffer}
\label{sec:l04-btb}

Predict not taken uses no information about the instruction being fetched.
Better predictors use its \emph{history}: what the instruction at this PC did
the previous times it was executed.\source{Lesson 4, videos 12--15; H\&P
§3.9; \cite{lee1984}.}  The next PC is predicted as a function of the PC and
of this history.  Because the instruction stored at an address does not
change, whether it is a branch and where a direct branch jumps to are the
same every time; and most branches go in the same direction most of the time.

\begin{definition}[Branch target buffer]
  A \term{branch target buffer}[分岐先バッファ]
  \index{BTB|see{branch target buffer}}(BTB) is a table, indexed by the PC
  of the instruction being fetched, that stores the predicted next PC for
  that instruction.
\end{definition}

In IF the BTB is read with the current PC, and its output is used as the
next PC.  When the branch resolves, the actual next PC is compared with the
prediction.  On a mismatch the wrongly fetched instructions are flushed and
the actual target is written into the BTB entry, so that the next execution
of the same instruction is predicted with the new target.

\subsection{A realistic BTB}

A BTB cannot have one entry for every possible instruction address.  It must
be read within the fetch cycle, in addition to the fetch itself, which
limits it to a small table---hundreds to a few thousand entries.\margin{Real
capacities: about 4096 BTB entries on Intel Skylake; AMD's Zen~2 uses two
levels, 512 entries in front of 7168 (vendor optimization guides).}  The PC is
therefore mapped to an entry by a simple hash: a few of its low-order bits.
Instructions executed close together in time, such as those of a loop, lie
close together in memory; they share their high-order bits but differ in the
low-order bits, which spreads them over different entries.  Two instructions
collide only if they are far apart in the code, and such instructions are
less likely to be active at the same time.

With fixed-length 4-byte instructions aligned to word boundaries, the two
lowest bits of every PC are zero and carry no information, so a BTB with
$2^k$ entries is indexed by bits $k+1$ to~2 (\cref{fig:l04-btb}).\tip{Entry
number $= (\text{PC} \gg 2) \mathbin{\&} (\text{entries} - 1)$: drop the two
always-zero bits, then keep the $k$ low bits.}  With
variable-length instructions (x86) the index starts at bit~0.

\begin{figure}[htbp]
  \centering
  % A realistic BTB: low-order PC bits (above the alignment bits) index a
% small table of predicted targets.
\begin{tikzpicture}[x=1cm,y=1cm, font=\footnotesize\sffamily, >={Stealth[length=1.8mm]}]
  % PC fields
  \node[anchor=east] at (-0.1,0.3) {PC};
  \draw[fill=white] (0,0) rectangle (3.0,0.6);   \node at (1.5,0.3) {31 \dots\ 10};
  \draw[fill=ln-accent!15] (3.0,0) rectangle (5.6,0.6); \node at (4.3,0.3) {9 \dots\ 2};
  \draw[fill=ln-box] (5.6,0) rectangle (6.4,0.6); \node at (6.0,0.3) {1\,0};
  \node[anchor=west, font=\scriptsize, text=ln-muted] at (6.5,0.3) {\iflangja{常に 00}{always 00}};
  % index
  \draw[->] (4.3,0) -- (4.3,-0.8) node[midway, right, font=\scriptsize] {\iflangja{インデックス（8 ビット）}{index (8 bits)}};
  % table: 6 rows of 0.45
  \fill[ln-accent!15] (3.0,-2.15) rectangle (5.6,-2.6);
  \foreach \i in {0,...,5} { \draw (3.0,-0.8-0.45*\i) rectangle (5.6,-1.25-0.45*\i); }
  \draw (3.7,-0.8) -- (3.7,-3.5);
  \foreach \i/\lab in {0/0,1/1,3/170,5/255} {
    \node[font=\scriptsize, text=ln-muted] at (3.35,-1.025-0.45*\i) {\lab}; }
  \foreach \i in {2,4} { \node[font=\scriptsize, text=ln-muted] at (3.35,-1.0-0.45*\i) {$\vdots$}; }
  \node at (4.65,-2.375) {\iflangja{分岐先}{target}};
  \node[font=\scriptsize, text=ln-muted, anchor=north] at (4.3,-3.55) {BTB};
  % lookup and update
  \draw[->, thick] (5.6,-2.375) -- (7.0,-2.375) node[right] {\iflangja{予測した次の PC}{predicted next PC}};
  \draw[->, dashed, ln-caution] (0.6,-2.375) -- (3.0,-2.375)
    node[pos=0.0, above right, font=\scriptsize, align=left] {\iflangja{分岐解決時に\\実際の分岐先を書く}{on resolution:\\write actual target}};
\end{tikzpicture}

  \caption{A 256-entry BTB indexed by PC bits 9--2.  The entry is read in IF
    and written with the actual target when a branch resolves.}
  \label{fig:l04-btb}
\end{figure}

\begin{example}
  \label{ex:l04-btb-index}
  In a 256-entry BTB, the instruction at \texttt{0x004012A8} uses bits 9--2 of
  its PC: $(\texttt{0x12A8} \gg 2) \mathbin{\&} \texttt{0xFF} =
  \texttt{0x4AA} \mathbin{\&} \texttt{0xFF} = \texttt{0xAA}$, i.e.\
  entry~170.  The next instruction, at \texttt{0x004012AC}, uses entry~171;
  the instruction at \texttt{0x004016A8}, 256 instructions further on, shares
  entry~170.
\end{example}

\section{Direction prediction}
\label{sec:l04-direction}

A BTB alone answers the direction question implicitly: an entry that holds
PC+4 means \enquote{not taken}.\source{Lesson 4, videos 16--21; H\&P §C.2.}
This spends a full target address on every non-branch and every not-taken
branch.  It is cheaper to split the prediction: a
\term{direction predictor}[方向予測器] decides whether the instruction is a
taken branch, and the BTB is consulted only when it says so.

\begin{definition}[Branch history table]
  A \term{branch history table}[分岐履歴テーブル]
  \index{BHT|see{branch history table}}(BHT) is a direction predictor in the
  form of a table indexed by low-order PC bits.  In its simplest form each entry is a single
  bit holding the outcome of the last execution of the branch: 0 selects
  PC+4 as the next PC, 1 selects the target from the BTB.  This is the
  \term{one-bit predictor}[1 ビット予測器]\index{1BP|see{one-bit predictor}}
  (1BP).
\end{definition}

When a branch resolves, its BHT entry is set to the actual outcome, and, if
the branch was taken, the target is written into the BTB
(\cref{fig:l04-bht-btb}).  This organization has three advantages.  The BHT
entry of a non-branch instruction is never written, stays~0, and correctly
predicts PC+4.\caution{Neither table is tagged: if a taken branch maps to the
same entry, that non-branch is predicted taken too, and the error is caught
only at decode.}  The BTB needs entries only for taken branches, which are far
fewer than all instructions.  And since a BHT entry is a single bit, the BHT
can have many more entries than the BTB at the same cost, so that different
branches rarely share an entry.

\begin{figure}[htbp]
  \centering
  % Fetch-stage predictor: the BHT predicts the direction, the BTB the target.
\begin{tikzpicture}[x=1cm,y=1cm, font=\footnotesize\sffamily, >={Stealth[length=1.8mm]},
  box/.style={draw, fill=ln-box, minimum height=0.6cm, inner sep=2pt}]
  \node[box, minimum width=0.9cm] (pc) at (0,0) {PC};
  % index bus
  \draw[->] (pc.east) -- (6.0,0) -- (6.0,-0.5);
  \draw[->] (1.2,0) -- (1.2,-0.9);
  \draw[->] (3.4,0) -- (3.4,-0.5);
  \fill (1.2,0) circle (1.2pt); \fill (3.4,0) circle (1.2pt);
  % +4
  \node[box, minimum width=0.8cm] (inc) at (1.2,-1.2) {$+4$};
  % BTB
  \draw[fill=ln-box] (2.5,-0.5) rectangle (4.3,-1.9);
  \foreach \y in {-0.78,-1.06,-1.34,-1.62} { \draw[ln-rule] (2.5,\y) -- (4.3,\y); }
  \node[anchor=east, font=\footnotesize\bfseries] at (2.45,-0.75) {BTB};
  % BHT
  \draw[fill=ln-box] (5.7,-0.5) rectangle (6.3,-1.9);
  \foreach \y in {-0.78,-1.06,-1.34,-1.62} { \draw[ln-rule] (5.7,\y) -- (6.3,\y); }
  \node[anchor=west, font=\footnotesize\bfseries] at (6.35,-0.75) {BHT};
  \node[anchor=west, font=\scriptsize, text=ln-muted] at (6.35,-1.15) {\iflangja{1 ビット/エントリ}{1 bit/entry}};
  % update
  \node[font=\scriptsize, text=ln-caution, align=center] (upd) at (5.0,-2.3) {\iflangja{EX で\\更新}{update\\from EX}};
  \draw[->, dashed, ln-caution] (upd.north) -- (4.3,-1.2);
  \draw[->, dashed, ln-caution] (upd.north) -- (5.7,-1.2);
  % mux
  \draw[fill=white] (0.7,-2.6) -- (3.9,-2.6) -- (3.3,-3.2) -- (1.3,-3.2) -- cycle;
  \draw[->] (inc.south) -- (1.2,-2.6);
  \draw[->] (3.4,-1.9) -- (3.4,-2.6);
  \node[font=\scriptsize] at (1.3,-2.8) {0};
  \node[font=\scriptsize] at (3.3,-2.8) {1};
  \draw[->] (6.0,-1.9) -- (6.0,-2.9) -- (3.6,-2.9)
    node[pos=0.72, below, font=\scriptsize] {\iflangja{成立?}{taken?}};
  % next PC feedback
  \draw[->] (2.3,-3.2) -- (2.3,-3.7) -- (-0.9,-3.7) -- (-0.9,0) -- (pc.west);
  \node[below, font=\scriptsize] at (0.7,-3.7) {\iflangja{次の PC}{next PC}};
\end{tikzpicture}

  \caption{Next-PC prediction with a BHT and a BTB.  The BHT bit selects
    between PC+4 and the BTB target; both tables are updated when the
    branch resolves.}
  \label{fig:l04-bht-btb}
\end{figure}

\begin{example}
  \label{ex:l04-1bp-trace}
  An inner loop processes three array elements each time it is entered:
  \begin{lstlisting}[language={[x86masm]Assembler}, numbers=none]
0x0400  Inner: LW   R2, 0(R3)     ; load element
0x0404         ADD  R5, R5, R2    ; add to sum
0x0408         ADDI R3, R3, 4     ; next element
0x040C         ADDI R1, R1, -1    ; decrement count
0x0410         BNE  R1, R0, Inner ; T, T, N per run
\end{lstlisting}
  With a 64-entry BHT (index bits 7--2) the five instructions use BHT
  entries 0 to~4; the \texttt{BNE} uses entry~4.  With an 8-entry BTB (bits
  4--2) the \texttt{BNE} uses BTB entry~4, which it shares with the
  instruction at \texttt{0x03F0}, eight instructions earlier.  If that
  instruction is not a branch, its BHT entry (entry~60) stays~0, so the BTB
  entry is used only by the branch.  \Cref{tab:l04-1bp-trace} follows the
  \texttt{BNE} over two runs of the loop, starting with all tables cleared.
\end{example}

\begin{table}[htbp]
  \centering\small
  \caption{One-bit prediction of the \texttt{BNE} of
    \cref{ex:l04-1bp-trace} (BHT entry~4) over two runs of the loop.}
  \label{tab:l04-1bp-trace}
  \begin{tabular}{@{}lcccccc@{}}
    \toprule
    Execution         & 1 & 2 & 3 & 4 & 5 & 6 \\
    \midrule
    BHT[4] before     & 0 & 1 & 1 & 0 & 1 & 1 \\
    Predicted next PC & \texttt{0x414} & \texttt{0x400} & \texttt{0x400} & \texttt{0x414} & \texttt{0x400} & \texttt{0x400} \\
    Outcome           & T & T & N & T & T & N \\
    Correct?          & no & yes & no & no & yes & no \\
    \bottomrule
  \end{tabular}
\end{table}

The first execution also writes \texttt{0x400} into BTB entry~4; from then
on the target is always right, but the direction is wrong in four of the six
executions.

\section{The two-bit predictor}
\label{sec:l04-2bp}

The one-bit predictor works well for branches that always go the same way,
but it has two weaknesses.\source{Lesson 4, videos 22--26; H\&P §C.2;
\cite{smith1981}.}
\begin{itemize}
  \item A branch that almost always goes one way suffers \emph{two}
    mispredictions per exception: one on the exception itself, and one on
    the next normal outcome, because the exception has flipped the bit.
  \item A branch whose direction changes often is predicted badly.  An
    alternating branch (T,~N,~T,~N,~\ldots) is mispredicted every time, and
    the short loop of \cref{tab:l04-1bp-trace} two times out of three.
\end{itemize}

\begin{definition}[Two-bit predictor]
  A \term{two-bit predictor}[2 ビット予測器]
  \index{2BP|see{two-bit predictor}}\index{2BC|see{two-bit predictor}}(2BP,
  also called a two-bit counter, 2BC) stores in each BHT
  entry a two-bit counter that is incremented when the branch is taken and
  decremented when it is not taken.
\end{definition}

The counter is a \term{saturating counter}[飽和カウンタ]: it stays at 11 when
incremented and at 00 when decremented.  Its more significant bit is the
\emph{prediction bit}, taken if~1.  The less significant bit is the
\emph{hysteresis bit}.  Together with the prediction bit it records how
strongly the counter believes its prediction: the counter is in a strong
state when the two bits are equal and in a weak state when they differ, so
the four states are strong not taken (00), weak not taken (01), weak taken
(10) and strong taken (11) (\cref{fig:l04-2bc}).\tip{Read the two bits as a
number 0--3: T adds one, N subtracts one, both clamped, and the prediction is
taken iff the value is at least~2.}  From a
strong state, two consecutive outcomes against the prediction are needed to
change it, so a single exception costs a single misprediction.

\begin{figure}[htbp]
  \centering
  % Two-bit saturating counter: T increments, N decrements.
\begin{tikzpicture}[font=\footnotesize\sffamily, >={Stealth[length=1.8mm]},
  s/.style={circle, draw, minimum size=0.95cm, inner sep=1pt},
  nt/.style={s, fill=ln-box}, tk/.style={s, fill=ln-accent!15}]
  \node[nt] (s0) at (0,0)   {00};
  \node[nt] (s1) at (2.5,0) {01};
  \node[tk] (s2) at (5.0,0) {10};
  \node[tk] (s3) at (7.5,0) {11};
  \draw[->] (s0) to[bend left=30] node[above] {T} (s1);
  \draw[->] (s1) to[bend left=30] node[above] {T} (s2);
  \draw[->] (s2) to[bend left=30] node[above] {T} (s3);
  \draw[->] (s3) to[bend left=30] node[below] {N} (s2);
  \draw[->] (s2) to[bend left=30] node[below] {N} (s1);
  \draw[->] (s1) to[bend left=30] node[below] {N} (s0);
  \draw[->] (s0) to[out=150,in=210,looseness=6] node[left] {N} (s0);
  \draw[->] (s3) to[out=30,in=-30,looseness=6] node[right] {T} (s3);
  \foreach \n/\lab in {s0/{\iflangja{強い不成立}{strong N}},s1/{\iflangja{弱い不成立}{weak N}},
                       s2/{\iflangja{弱い成立}{weak T}},s3/{\iflangja{強い成立}{strong T}}} {
    \node[font=\scriptsize, text=ln-muted, below=5mm of \n] {\lab};
  }
  \draw[ln-rule] (-0.6,-1.35) -- (3.1,-1.35);
  \draw[ln-rule] (4.4,-1.35) -- (8.1,-1.35);
  \node[font=\scriptsize, below] at (1.25,-1.35) {\iflangja{予測: 不成立}{predict not taken}};
  \node[font=\scriptsize, below] at (6.25,-1.35) {\iflangja{予測: 成立}{predict taken}};
\end{tikzpicture}

  \caption{State diagram of the two-bit saturating counter.  T (taken)
    moves right, N (not taken) moves left; the upper bit is the prediction.}
  \label{fig:l04-2bc}
\end{figure}

\begin{table}[htbp]
  \centering\small
  \caption{The loop branch of \cref{ex:l04-1bp-trace} over three runs of the
    loop with a two-bit counter starting at 00.  Mispredictions are marked
    $\times$.}
  \label{tab:l04-2bc-trace}
  \setlength{\tabcolsep}{4.5pt}
  \begin{tabular}{@{}lccccccccc@{}}
    \toprule
    Execution       & 1 & 2 & 3 & 4 & 5 & 6 & 7 & 8 & 9 \\
    \midrule
    Outcome         & T & T & N & T & T & N & T & T & N \\
    Counter before  & 00 & 01 & 10 & 01 & 10 & 11 & 10 & 11 & 11 \\
    Prediction      & N & N & T & N & T & T & T & T & T \\
    2BP             & $\times$ & $\times$ & $\times$ & $\times$ & & $\times$ & & & $\times$ \\
    1BP (for comparison) & $\times$ & & $\times$ & $\times$ & & $\times$ & $\times$ & & $\times$ \\
    \bottomrule
  \end{tabular}
\end{table}

\Cref{tab:l04-2bc-trace} shows the effect on the short loop.  Once the
counter has reached strong taken (11, before execution~6), a single loop
exit only moves it down to weak taken (10) and the next taken outcome
restores it, so only the exit is mispredicted: one misprediction per run of
the loop instead of two.\caution{From a \emph{weak} state one contrary
outcome is enough: a two-bit counter does not always need two mispredictions
to change its prediction.}

\subsection{Initialization and more bits}

A counter initialized to a strong state costs two mispredictions if the
branch actually goes the other way; one initialized to a weak state costs
only one.  The weak states, however, are also where the counter is most
vulnerable: an alternating branch T,~N,~T,~N,~\ldots{} that starts in 01
moves between 01 and 10 and is mispredicted every time, whereas starting in
00 it is mispredicted only on every other execution.  Neither choice is
best for every branch: each initial state has outcome sequences for which it
does worse.  Counters are therefore simply reset to the simplest state,
00.\margin{No
predictor is immune: for any predictor there is an outcome sequence that
always goes against its current prediction.}

Counters with three or more bits add hysteresis: more consecutive exceptions
are needed to change the prediction.  This helps only for branches whose
exceptions come in runs, it makes the counter slower to follow a branch that
really changes its behaviour, and it costs more bits.  The step from one to
two bits is a large improvement; further bits gain little, which is why the
two-bit counter is the standard building block.

\section{History-based predictors}
\label{sec:l04-history}

Some branches follow a regular pattern that counters cannot capture, because
counters only track which direction is more frequent.\source{Lesson 4, videos
27--33; H\&P §3.3; \cite{yeh1991}.}  A branch that tests whether a loop index
is even alternates T,~N,~T,~N; the loop branch of \cref{ex:l04-1bp-trace}
repeats T,~T,~N.  Both are perfectly predictable once the most recent
outcomes are known.

\begin{definition}[History-based predictor]
  A \term{history-based predictor}[履歴ベース予測器] records the most recent
  outcomes of a branch (its history) and uses them to select one of several
  counters, each of which learns what the branch does after one particular
  history.
\end{definition}

\subsection{One-bit history}

The simplest version stores in each BHT entry one history bit, the last
outcome of the branch, and two two-bit counters: one used when the last
outcome was N, one when it was T.  To predict, the history bit selects a
counter and the counter's upper bit is the prediction.  To update, the
selected counter is incremented or decremented with the actual outcome, and
the outcome becomes the new history bit.  An entry needs $1 + 2\times 2 = 5$
bits.

For the alternating branch, the counter selected after N sees only T and the
counter selected after T sees only N.  After each has been trained, the
branch is predicted correctly every time.  For the pattern T,~T,~N this is not
enough: after a T the branch is sometimes taken and sometimes not, so the
counter selected after T keeps mispredicting.

\subsection{Two-bit and \texorpdfstring{$N$}{N}-bit history}

With two history bits and four counters (10 bits per entry,
\cref{fig:l04-history}), each of the four possible histories NN, NT, TN and
TT has its own counter.  For T,~T,~N the history TT is always followed by N,
while TN and NT are always followed by T.  \Cref{tab:l04-history-trace}
traces the loop branch: after the counters have been trained, every
execution is predicted correctly.  The counter for NN is never used again.

\begin{figure}[htbp]
  \centering
  % Entry of a two-bit history predictor: the history selects one of four
% two-bit counters.
\begin{tikzpicture}[x=1cm,y=1cm, font=\footnotesize\sffamily, >={Stealth[length=1.8mm]}]
  \node[draw, fill=ln-box, minimum width=0.9cm, minimum height=0.6cm] (pc) at (-1.7,0) {PC};
  \draw[->] (pc.east) -- (0,0) node[midway, above, font=\scriptsize] {\iflangja{索引}{index}};
  \draw[fill=ln-accent!15] (0,-0.3) rectangle (1.4,0.3); \node at (0.7,0) {$h_1h_0$};
  \foreach \i/\lab in {0/NN,1/NT,2/TN,3/TT} {
    \draw[fill=ln-box] (1.4+1.25*\i,-0.3) rectangle ++(1.25,0.6);
    \node at (2.025+1.25*\i,0) {$C_{\mathrm{\lab}}$};
    \draw[->] (2.025+1.25*\i,-0.3) -- (2.025+1.25*\i,-1.1);
  }
  \node[font=\scriptsize, text=ln-muted, anchor=south] at (0.7,0.3) {\iflangja{履歴 2 ビット}{2-bit history}};
  \node[font=\scriptsize, text=ln-muted, anchor=south] at (3.9,0.3) {\iflangja{2 ビットカウンタ $\times 4$}{four 2-bit counters}};
  % 4-to-1 mux
  \draw[fill=white] (1.6,-1.1) -- (6.2,-1.1) -- (4.8,-1.7) -- (3.0,-1.7) -- cycle;
  \draw[->] (0.7,-0.3) -- (0.7,-1.4) -- (2.3,-1.4);
  \node[font=\scriptsize, below] at (1.3,-1.4) {\iflangja{選択}{select}};
  \draw[->] (3.9,-1.7) -- (3.9,-2.2) node[below, font=\scriptsize]
    {\iflangja{予測 = 選ばれたカウンタの上位ビット}{prediction = MSB of the selected counter}};
\end{tikzpicture}

  \caption{An entry of a two-bit history predictor.  The history bits select
    one of four two-bit counters; after the branch resolves, the selected
    counter is updated and the outcome is shifted into the history.}
  \label{fig:l04-history}
\end{figure}

\begin{table}[htbp]
  \centering\small
  \caption{Two-bit history predictor on the loop branch of
    \cref{ex:l04-1bp-trace}, starting with history NN and all counters at
    00.  The history is written oldest outcome first; mispredictions are
    marked $\times$.}
  \label{tab:l04-history-trace}
  \setlength{\tabcolsep}{3.6pt}
  \begin{tabular}{@{}lcccccccccccc@{}}
    \toprule
    Execution      & 1 & 2 & 3 & 4 & 5 & 6 & 7 & 8 & 9 & 10 & 11 & 12 \\
    \midrule
    History        & NN & NT & TT & TN & NT & TT & TN & NT & TT & TN & NT & TT \\
    Counter used   & 00 & 00 & 00 & 00 & 01 & 00 & 01 & 10 & 00 & 10 & 11 & 00 \\
    Prediction     & N & N & N & N & N & N & N & T & N & T & T & N \\
    Outcome        & T & T & N & T & T & N & T & T & N & T & T & N \\
    Mispredicted   & $\times$ & $\times$ & & $\times$ & $\times$ & & $\times$ & & & & & \\
    \bottomrule
  \end{tabular}
\end{table}

\begin{definition}[$N$-bit history predictor]
  An \term[N-bit history predictor]{$N$-bit history predictor}[N ビット履歴予測器]
  keeps the last $N$ outcomes of each branch and $2^N$ two-bit counters per
  entry, one for each possible history.  An entry costs
  \begin{equation}
    N + 2 \cdot 2^N \ \text{bits}.
    \label{eq:l04-history-cost}
  \end{equation}
\end{definition}

After training, an $N$-bit history predictor predicts every repeating
pattern whose period is at most $N+1$, because within such a pattern the
last $N$ outcomes determine the next one.  The price grows exponentially
with~$N$, and most of it is wasted: a branch with period $p$ uses at most $p$
of its $2^N$ counters.  Longer histories also take longer to train, since
every counter that is used has to be trained separately.

\begin{example}
  \label{ex:l04-history-cost}
  A loop that always runs 10 iterations has a branch with the pattern
  T\,T\,T\,T\,T\,T\,T\,T\,T\,N, of period~10.  Predicting it perfectly requires
  $N = 9$, i.e.\ $9 + 2\cdot 512 = 1033$ bits per BHT entry, of which the
  branch uses 10 of the 512 counters.  A BHT of 1024 such entries takes more
  than a megabit.
\end{example}

\begin{keypoint}
  History lets a predictor learn patterns, not only biases.  Giving every
  branch its own set of $2^N$ counters, however, makes long histories
  prohibitively expensive, and nearly all of those counters are never used.
\end{keypoint}

\section{Shared counters: PShare and GShare}
\label{sec:l04-share}

The waste disappears if the counters are not private to a
branch.\source{Lesson 4, videos 34--39; H\&P §3.3; \cite{mcfarling1993}.}
The BHT keeps only an $N$-bit history per entry, and all branches share one
large table of two-bit counters.

\begin{definition}[Pattern history table]
  A \term{pattern history table}[パターン履歴テーブル]
  \index{PHT|see{pattern history table}}(PHT) is a table of two-bit counters shared by all
  branches and indexed by a combination of the branch's history and its PC,
  typically their bitwise XOR.
\end{definition}

A branch now occupies only as many counters as it has distinct histories.
An always-taken branch has the history T\,T\,\ldots\,T and uses a single
counter, an alternating branch uses two, and the 10-iteration loop branch of
\cref{ex:l04-history-cost} uses ten.  Mixing in the PC bits sends different
branches with the same history to different counters.  Collisions remain
possible, but with a large PHT they are rare.

\begin{example}
  A BHT of 1024 entries holding 12-bit histories takes \num{12288} bits, and
  a PHT of $2^{12} = 4096$ counters takes \num{8192} bits: \num{20480} bits
  in total.  A 1024-entry predictor with 12-bit history and private counters
  would need $1024 \times (12 + 2\cdot 4096) = \num{8400896}$ bits, about
  410 times as much.
\end{example}

\subsection{Private and global history}

The history that indexes the PHT can be chosen in two ways
(\cref{fig:l04-pshare-gshare}).

\begin{definition}[PShare]
  The \term{private branch history}[ローカル分岐履歴] of a branch is the
  sequence of its own most recent outcomes, kept in its BHT entry.  A
  \term{PShare predictor}[PShare 予測器] indexes the shared PHT with the
  private history XORed with the PC.
\end{definition}

\begin{definition}[GShare]
  The \term{global branch history}[グローバル分岐履歴] is the sequence of the
  most recent outcomes of all branches, kept in a single shift register into
  which every resolved branch shifts its outcome.  A \term{GShare
  predictor}[GShare 予測器] indexes the shared PHT with the global history
  XORed with the PC.
\end{definition}

\begin{figure}[htbp]
  \centering
  % PShare (private history) and GShare (global history), both with a shared
% pattern history table of two-bit counters.
\begin{tikzpicture}[x=1cm,y=1cm, font=\footnotesize\sffamily, >={Stealth[length=1.8mm]},
  box/.style={draw, fill=ln-box, minimum height=0.6cm, minimum width=0.9cm, inner sep=2pt},
  xor/.style={inner sep=0pt, font=\large}]
  % ---------------- (a) PShare
  \node[anchor=west, font=\footnotesize\bfseries] at (-0.45,1.45) {(a) PShare};
  \node[box] (pca) at (0,0) {PC};
  \draw[fill=ln-box] (1.3,-0.7) rectangle (2.9,0.7);
  \foreach \y in {-0.35,0.35} { \draw[ln-rule] (1.3,\y) -- (2.9,\y); }
  \fill[ln-accent!15] (1.31,-0.34) rectangle (2.89,0.34);
  \node[font=\scriptsize, text=ln-muted] at (2.1,0.95) {\iflangja{BHT（ローカル履歴）}{BHT (private histories)}};
  \draw[->] (pca.east) -- (1.3,0);
  \node[xor] (xa) at (4.0,0) {$\oplus$};
  \draw[->] (2.9,0) -- (xa.west);
  \draw[->] (pca.south) -- (0,-1.0) -- (4.0,-1.0) -- (xa.south);
  \draw[fill=ln-box] (5.0,-0.7) rectangle (6.2,0.7);
  \foreach \y in {-0.35,0,0.35} { \draw[ln-rule] (5.0,\y) -- (6.2,\y); }
  \node[font=\scriptsize, text=ln-muted] at (5.6,0.95) {\iflangja{PHT（2 ビットカウンタ）}{PHT (2-bit counters)}};
  \draw[->] (xa.east) -- (5.0,0);
  \draw[->] (6.2,0) -- (7.0,0) node[right] {\iflangja{予測}{prediction}};
  % ---------------- (b) GShare
  \begin{scope}[yshift=-3.4cm]
  \node[anchor=west, font=\footnotesize\bfseries] at (-0.45,1.45) {(b) GShare};
  \foreach \i/\b in {0/1,1/0,2/1,3/1,4/0,5/1} {
    \draw[fill=ln-accent!15] (1.3+0.2667*\i,-0.3) rectangle ++(0.2667,0.6);
    \node[font=\scriptsize\ttfamily] at (1.4333+0.2667*\i,0) {\b};
  }
  \node[font=\scriptsize, text=ln-muted, align=center] at (2.1,0.8) {\iflangja{グローバル履歴\\（シフトレジスタ）}{global history\\(shift register)}};
  \node[xor] (xb) at (4.0,0) {$\oplus$};
  \draw[->] (2.9,0) -- (xb.west);
  \node[box] (pcb) at (0,-1.0) {PC};
  \draw[->] (pcb.east) -- (4.0,-1.0) -- (xb.south);
  \draw[fill=ln-box] (5.0,-0.7) rectangle (6.2,0.7);
  \foreach \y in {-0.35,0,0.35} { \draw[ln-rule] (5.0,\y) -- (6.2,\y); }
  \node[font=\scriptsize, text=ln-muted] at (5.6,0.95) {\iflangja{PHT（2 ビットカウンタ）}{PHT (2-bit counters)}};
  \draw[->] (xb.east) -- (5.0,0);
  \draw[->] (6.2,0) -- (7.0,0) node[right] {\iflangja{予測}{prediction}};
  \end{scope}
\end{tikzpicture}

  \caption{Predictors with shared counters.  (a) PShare combines the
    branch's own history with the PC; (b) GShare combines the global history
    of all branches with the PC.}
  \label{fig:l04-pshare-gshare}
\end{figure}

PShare is suited to branches with patterns of their own: alternating
branches and short loops.  GShare is suited to \emph{correlated} branches,
whose outcome depends on the outcomes of other branches executed just before.

\needspace{9\baselineskip}
\begin{example}
  In the fragment below the second branch is taken exactly when the first is
  not.
  \begin{lstlisting}[language={[x86masm]Assembler}, numbers=none]
        BNE  R1, R0, Skip1   ; skip unless R1 == 0
        ADDI R2, R2, 1       ; count zeros
Skip1:  BEQ  R1, R0, Skip2   ; skip if R1 == 0
        ADD  R3, R3, R1      ; sum nonzero values
Skip2:  ...
\end{lstlisting}
  If the value in R1 depends on the data, the private history of either
  branch looks random, and PShare predicts it poorly.\margin{A branch whose
  outcome depends on data may be cheaper to remove than to predict: Lesson 5
  replaces such branches with predication.}  The global history
  seen by the second branch contains the outcome of the first, so once the
  PHT counters for these histories have been trained, GShare predicts the
  second branch perfectly, barring aliasing in the PHT.
\end{example}

Neither is better for all branches.\caution{Only GShare is a standard name
(McFarling): what the lecture calls PShare is a local- or private-history
predictor in the literature, as in Yeh and Patt's PAs schemes.}  GShare can in
principle learn a branch's
own pattern as well, but the outcomes of the other branches executed in
between are interleaved in the global history, so the same pattern needs a
considerably longer history---and a larger PHT---than in PShare.  Since
programs contain both kinds of branches, the natural next step is to use both
predictors.

\section{Tournament and hierarchical predictors}
\label{sec:l04-hybrid}

Since different branches are best served by different predictors, a
processor can combine several predictors and use, for each branch, the one
that suits it.\source{Lesson 4, videos 40--43; H\&P §3.3;
\cite{mcfarling1993}.}  The combination can be symmetric or layered.

\subsection{Tournament predictors}

A \term{tournament predictor}[トーナメント予測器] runs two predictors, for
example GShare and PShare, side by side and decides separately for each
branch which of them to trust.

\begin{definition}[Meta-predictor]
  The \term{meta-predictor}[メタ予測器] of a tournament predictor is a
  table of two-bit counters, indexed by the PC, whose counter for a branch
  predicts which of the two component predictors will be right
  (\cref{fig:l04-tournament}).
\end{definition}

Both component predictors are trained on every branch as usual.\margin{The
Alpha 21264 combined 1024 local histories of 10 bits, a 4096-entry global
table indexed by 12 history bits and a 4096-entry chooser: about 29 kilobits
(H\&P §3.3).}  The meta-predictor counter changes only when the components
disagree in correctness: it moves towards the component that was right.
When both are right or both are wrong, it stays unchanged.

\begin{figure}[htbp]
  \centering
  % Tournament predictor: a meta-predictor chooses between two predictors.
\begin{tikzpicture}[x=1cm,y=1cm, font=\footnotesize\sffamily, >={Stealth[length=1.8mm]},
  box/.style={draw, fill=ln-box, minimum height=0.7cm, minimum width=2.5cm, inner sep=2pt, align=center}]
  \node[draw, fill=ln-box, minimum width=0.9cm, minimum height=0.6cm] (pc) at (0,0) {PC};
  \node[box] (p1) at (3.0,1.2) {PShare};
  \node[box] (p2) at (3.0,0) {GShare};
  \node[box, fill=ln-accent!15] (meta) at (3.0,-1.3) {\iflangja{メタ予測器}{meta-predictor}\\[-1pt]{\scriptsize\iflangja{（2 ビットカウンタ）}{(2-bit counters)}}};
  \draw (pc.east) -- (1.2,0);
  \draw (1.2,1.2) -- (1.2,-1.3);
  \fill (1.2,0) circle (1.2pt);
  \draw[->] (1.2,1.2) -- (p1.west);
  \draw[->] (1.2,0) -- (p2.west);
  \draw[->] (1.2,-1.3) -- (meta.west);
  % mux
  \draw[fill=white] (5.4,1.6) -- (6.0,1.2) -- (6.0,-0.4) -- (5.4,-0.8) -- cycle;
  \draw[->] (p1.east) -- (5.4,1.2);
  \draw[->] (p2.east) -- (5.4,0);
  \draw[->] (meta.east) -- (5.7,-1.3) -- (5.7,-0.6);
  \node[font=\scriptsize, below] at (5.0,-1.3) {\iflangja{どちら?}{which one?}};
  \draw[->] (6.0,0.4) -- (6.8,0.4) node[right] {\iflangja{予測}{prediction}};
\end{tikzpicture}

  \caption{A tournament predictor.  The meta-predictor selects, per branch,
    the component predictor whose prediction is used.}
  \label{fig:l04-tournament}
\end{figure}

\subsection{Hierarchical predictors}

A tournament predictor combines two good---and therefore expensive---predictors
and updates both on every branch.  Most branches, however, are easy to
predict.

\begin{definition}[Hierarchical predictor]
  A \term{hierarchical predictor}[階層型予測器] combines a cheap predictor,
  which handles all branches, with one or more better and more expensive
  predictors that are used and trained only for the branches the cheap
  predictor handles poorly.
\end{definition}

Because the expensive predictor has far fewer branches to learn, it can be
much smaller than in a tournament predictor for the same accuracy, or more
accurate for the same size.\margin{AMD's Zen~2 layers by speed instead: a fast
perceptron predicts, a slower TAGE predictor overrides it---about
\SI{30}{\percent} fewer mispredictions than Zen~1 (AMD).}

A typical arrangement has three levels (\cref{fig:l04-hierarchical}).  A large
table of two-bit counters indexed by the PC predicts every branch.  Above it,
a small local-history predictor and a small global-history predictor hold
entries only for selected branches.  Their entries are \emph{tagged} with
bits of the PC, so that a lookup can tell whether the branch has an entry.
The prediction comes from the highest level with a matching entry.  The
counter table is always updated, and a local or global entry that matches
the branch is updated as well; a branch that the counter table mispredicts
is given an entry in the local predictor, and a branch that the local
predictor mispredicts is given one in the global predictor.  Branches that
the counters predict well take entries in the expensive levels only rarely,
so those small tables are left for the branches that need
them.\sidenote{Intel's Pentium~M (2003) is the classic example: a table of
two-bit counters covering all branches, backed by tagged local-history and
global-history predictors, a loop detector that learns how many iterations a
loop branch runs, and a separate predictor for indirect branches (Gochman et
al., \emph{Intel Technology Journal}~7(2), 2003).}

\begin{figure}[htbp]
  \centering
  % Hierarchical predictor: tagged upper levels override a cheap 2BC table.
\begin{tikzpicture}[x=1cm,y=1cm, font=\footnotesize\sffamily, >={Stealth[length=1.8mm]},
  box/.style={draw, fill=ln-box, minimum height=0.8cm, minimum width=3.0cm, inner sep=2pt, align=center}]
  \node[draw, fill=ln-box, minimum width=0.9cm, minimum height=0.6cm] (pc) at (0,0) {PC};
  \node[box, fill=ln-accent!15] (g) at (3.2,1.3) {\iflangja{グローバル履歴予測器}{global-history predictor}\\[-1pt]{\scriptsize\iflangja{小・タグ付き}{small, tagged}}};
  \node[box, fill=ln-accent!8] (l) at (3.2,0) {\iflangja{ローカル履歴予測器}{local-history predictor}\\[-1pt]{\scriptsize\iflangja{小・タグ付き}{small, tagged}}};
  \node[box] (b) at (3.2,-1.3) {\iflangja{2 ビットカウンタテーブル}{2-bit counter table}\\[-1pt]{\scriptsize\iflangja{大・タグなし}{large, untagged}}};
  \draw (pc.east) -- (1.2,0);
  \draw (1.2,1.3) -- (1.2,-1.3);
  \fill (1.2,0) circle (1.2pt);
  \foreach \n in {g,l,b} { \draw[->] (1.2,0 |- \n.west) -- (\n.west); }
  \node[draw, fill=white, align=left, font=\scriptsize, minimum height=3.4cm, text width=2.3cm] (sel) at (6.8,0)
    {\iflangja{1.\ グローバルがヒット\\\quad$\to$ それを使う\\[3pt]2.\ ローカルがヒット\\\quad$\to$ それを使う\\[3pt]3.\ それ以外\\\quad$\to$ 2BC を使う}{1.\ global hit\\\quad$\to$ use global\\[3pt]2.\ local hit\\\quad$\to$ use local\\[3pt]3.\ otherwise\\\quad$\to$ use 2BC}};
  \foreach \n in {g,l,b} { \draw[->] (\n.east) -- (sel.west |- \n.east); }
  \draw[->] (sel.east) -- ++(0.6,0) node[right] {\iflangja{予測}{prediction}};
\end{tikzpicture}

  \caption{A hierarchical predictor.  Tagged upper levels override the
    two-bit counter table for the branches they hold.}
  \label{fig:l04-hierarchical}
\end{figure}

\Cref{tab:l04-which} summarizes which kind of predictor handles which kind of
branch.

\begin{table}[htbp]
  \centering\small
  \caption{Branch behaviour and the simplest predictor that handles it well.}
  \label{tab:l04-which}
  \begin{tabular}{@{}p{0.52\linewidth}p{0.4\linewidth}@{}}
    \toprule
    Branch behaviour & Predictor \\
    \midrule
    Almost always the same direction & Two-bit counter \\
    Own repeating pattern (alternation, short loop) & Private history (PShare) \\
    Correlated with preceding branches & Global history (GShare) \\
    Mixture of the above in one program & Tournament or hierarchical \\
    Function return & Return address stack (next section) \\
    \bottomrule
  \end{tabular}
\end{table}

\section{The return address stack}
\label{sec:l04-ras}

A function return (\texttt{RET}) is an unconditional branch, so its direction
is trivially predicted as taken.\source{Lesson 4, videos 44--47; H\&P §3.9;
\cite{kaeli1991}.}  Its target, however, depends on where the function was
called from.  The BTB remembers only the most recent target, so when a
function is called from different places in turn, every return is
mispredicted.

\begin{example}
  A function is called from two places:
  \begin{lstlisting}[language={[x86masm]Assembler}, numbers=none]
0x1000         CALL Func    ; return address 0x1004
  ...
0x1040         CALL Func    ; return address 0x1044
  ...
0x2000  Func:  ...
0x2030         RET          ; to 0x1004, then to 0x1044
\end{lstlisting}
  When the first return resolves, the BTB entry of the \texttt{RET}
  receives \texttt{0x1004}; that return is itself mispredicted, since the
  entry was still empty.  The second return, which goes to \texttt{0x1044},
  is mispredicted and overwrites the entry.  If calls from the two places
  alternate, every return is mispredicted.
\end{example}

\begin{definition}[Return address stack]
  A \term{return address stack}[リターンアドレススタック]
  \index{RAS|see{return address stack}}(RAS) is a small hardware stack in the fetch stage.  When a
  call is fetched, its return address (PC+4) is pushed; when a return is
  fetched, the top entry is popped and used as the predicted target.
\end{definition}

Since calls and returns nest, the RAS predicts returns almost perfectly
(\cref{fig:l04-ras}).  Being small, it can fill up when calls nest
deeply.\margin{Real return stacks hold 16 entries per hardware thread on
Intel cores and 32 on AMD's Zen~2 (vendor documentation).}
When a call finds the RAS full, the push wraps around and overwrites the
oldest entry.  Returns occur in reverse order of the calls, so the entries
needed soonest are the most recent ones, and most calls and returns happen
deep in the call tree.  Overwriting the oldest entry loses only the
predictions for returns from the outermost functions, which are rare, whereas
discarding the new return address would lose the predictions for the
frequent calls at the deepest level.

\begin{figure}[htbp]
  \centering
  % Return address stack: calls push, returns pop; a full stack wraps around.
\begin{tikzpicture}[x=1cm,y=1cm, font=\footnotesize\sffamily, >={Stealth[length=1.8mm]}]
  \foreach \i/\v in {0/0x1004,1/0x2210,2/0x2A08,3/0x3104} {
    \draw[fill=ln-box] (0,0.55*\i) rectangle (1.8,0.55*\i+0.55);
    \node[font=\footnotesize\ttfamily] at (0.9,0.55*\i+0.275) {\v};
  }
  \draw[very thick, ln-accent] (0,1.65) rectangle (1.8,2.2);
  \node[font=\scriptsize, text=ln-muted, anchor=north] at (0.9,-0.05) {\iflangja{最も古い}{oldest}};
  % pop from the top
  \draw[->, thick] (1.85,1.925) -- (3.0,1.925);
  \node[anchor=west, font=\scriptsize] at (3.05,1.925) {\iflangja{RET: 先頭をポップ → 予測した分岐先}{RET: pop top $\to$ predicted target}};
  % push onto a full stack wraps around to the oldest entry
  \draw[thick] (3.0,2.6) -- (0.9,2.6);
  \draw[->, thick, dashed] (0.9,2.6) -- (-0.4,2.6) to[out=180,in=180,looseness=1.6] (-0.05,0.275);
  \node[anchor=west, font=\scriptsize] at (3.05,2.6) {\iflangja{CALL: 戻りアドレス（PC+4）をプッシュ}{CALL: push return address (PC+4)}};
  \node[anchor=east, align=right, font=\scriptsize, text=ln-muted] at (-1.45,1.35)
    {\iflangja{満杯のとき:\\プッシュで\\最古を上書き}{when full:\\push overwrites\\oldest}};
\end{tikzpicture}

  \caption{A four-entry return address stack.  Calls push, returns pop; when
    the stack is full, a push wraps around and overwrites the oldest entry.}
  \label{fig:l04-ras}
\end{figure}

\subsection{Recognizing a return}

The RAS has to be used while the \texttt{RET} is being fetched, before it has
been decoded.  There are two ways to know that the instruction is a return.
\begin{itemize}
  \item \textbf{Predict it.}  The predictor tables store, per PC, whether the
    instruction is a return; the entry is set the first time the instruction
    is decoded.  Since the instruction at a given address does not change,
    this prediction is almost always right.
  % keeps this item with the lesson summary, which alone would fill a page
  \needspace{9\baselineskip}
  \item \textbf{Predecode it.}  With \term{predecoding}[プリデコード], the
    processor partially decodes each instruction when it is brought from
    main memory into the processor's cache, and stores a few extra bits with
    the cached instruction, for example whether it is a branch or a return.  The fetch
    stage reads these bits together with the instruction and knows its kind
    immediately.
\end{itemize}

\begin{keypoint}[Lesson summary]
  Every fetch is a prediction.  A misprediction wastes the work of every
  instruction it flushes (\cref{eq:l04-relative}), and the number of flushed
  instructions grows with pipeline depth and width.  Predict not taken is
  free but inaccurate.  A BTB predicts
  targets and a BHT predicts directions, best with two-bit saturating
  counters.  History-based predictors learn patterns; sharing their counters
  in a PHT makes long histories affordable, with private history (PShare)
  for self-repeating branches and global history (GShare) for correlated
  ones.  Tournament and hierarchical predictors combine several predictors,
  and a return address stack predicts function returns.
\end{keypoint}

\end{document}
